\chapter{Optimization Problems in Economics}

Economics is the science of choice under scarcity, and its mathematical
core is the optimization problem: an agent
selects the values of some variables under their control so as to make an
objective as large (or as small) as possible, while respecting the
limits imposed by budgets, technology, and time. A consumer maximises
utility subject to a budget; a firm minimises the cost of producing a
given output, or maximises profit outright; an investor trades off risk
against return. Before any calculus is done, every such problem must be
\emph{formulated}: what is chosen, what is being optimised, and what
constrains the choice. This chapter drills exactly that skill, together
with the vocabulary (objective, decision variables, constraints,
feasible set, local versus global optima) and the question of whether an
optimum exists at all.

As a warm-up for the chapters that follow, we also practise
\emph{first-order thinking} in one variable: setting a derivative to zero
to locate candidates and using the second derivative to classify them.
Many economic problems with one constraint reduce to this case by
substituting the constraint into the objective; the full multivariate and
Lagrangian machinery is developed later in the book.

\section{Core Rules}

\begin{rulebox}{The general optimization problem}
An optimization problem has the standard form
\[
  \max_{x \in S} f(x)
  \qquad\text{or}\qquad
  \min_{x \in S} f(x),
\]
where
\begin{itemize}[leftmargin=1.6em,itemsep=0.3ex]
  \item $x = (x_1,\dots,x_n)$ are the \textbf{decision variables}: the quantities the agent controls;
  \item $f : \R^n \to \R$ is the \textbf{objective function}: the quantity being maximised or minimised;
  \item $S \subseteq \R^n$ is the \textbf{feasible set}: the choices allowed by the
        \textbf{constraints}, typically written
        as equalities $g(x)=b$ or inequalities $g(x)\le b$, plus sign
        restrictions such as $x \ge 0$.
\end{itemize}
A point $x \in S$ is called \emph{feasible}; a feasible point that
optimises $f$ over $S$ is a \emph{solution} or \emph{optimum}, written
$x^{*}$, with \emph{optimal value} $f(x^{*})$.
\end{rulebox}

\begin{rulebox}{Maximisation and minimisation are one theory}
For any $f$ and $S$,
\[
  \max_{x \in S} f(x) \;=\; -\min_{x \in S}\bigl(-f(x)\bigr),
\]
and both problems have exactly the same solutions $x^{*}$. Every result
stated for minima therefore has a mirror image for maxima (and vice
versa): flip the sign of the objective, not the theory.
\end{rulebox}

\begin{rulebox}{Local versus global optima}
Let $x^{*} \in S$.
\begin{itemize}[leftmargin=1.6em,itemsep=0.3ex]
  \item $x^{*}$ is a \textbf{global maximum} if
        $f(x^{*}) \ge f(x)$ for \emph{all} $x \in S$;
  \item $x^{*}$ is a \textbf{local maximum} if
        $f(x^{*}) \ge f(x)$ for all feasible $x$ in some neighbourhood of
        $x^{*}$.
\end{itemize}
Global minima and local minima are defined with $\le$. Every global
optimum is a local optimum; the converse is false in general. Calculus
conditions (derivatives set to zero) detect \emph{local} candidates only;
promoting a local result to a global one needs an extra argument,
concavity/convexity of $f$, or a direct comparison of all candidates.
\end{rulebox}

\begin{rulebox}{Existence: the Weierstrass theorem (informal)}
If the feasible set $S$ is nonempty, \emph{closed} (contains its boundary
points) and \emph{bounded} (fits inside some ball) (in short,
\emph{compact}) and the objective $f$ is
\emph{continuous} on $S$, then the problem has at least one global
maximum and at least one global minimum on $S$. Both hypotheses matter:
\begin{itemize}[leftmargin=1.6em,itemsep=0.3ex]
  \item $\max x$ over the open interval $(0,1)$ has no solution (the
        supremum $1$ is never attained ($S$ not closed);
  \item $\max x$ over $[0,\infty)$ has no solution) the objective is
        unbounded ($S$ not bounded).
\end{itemize}
The theorem is sufficient, not necessary: optima can exist on non-compact
sets (e.g.\ $x + 1/x$ on $(0,\infty)$ attains its minimum at $x=1$), but
then existence must be argued separately.
\end{rulebox}

\begin{rulebox}{Constrained versus unconstrained}
A problem is \textbf{unconstrained} if the feasible
set is all of $\R^n$ or an open set: the optimum, if it exists, is an
interior point where derivatives vanish. It is \textbf{constrained} if constraints limit the choice. Two practical
remarks:
\begin{itemize}[leftmargin=1.6em,itemsep=0.3ex]
  \item A constraint that is \emph{slack} (not binding) at the optimum can
        be dropped without changing the solution: a problem over $q \ge 0$
        whose solution satisfies $q^{*} > 0$ is \emph{effectively}
        unconstrained.
  \item A single equality constraint can often be \emph{substituted} into
        the objective, converting a constrained problem into an
        unconstrained one with fewer variables. This is the main
        method of this chapter.
\end{itemize}
\end{rulebox}

\begin{rulebox}{The economic archetypes}
Four problem families recur throughout economics:
\begin{itemize}[leftmargin=1.6em,itemsep=0.3ex]
  \item \textbf{Utility maximisation}: choose a
        consumption bundle $(x,y)\ge 0$ to
        \[
          \max\; U(x,y) \quad\text{subject to}\quad p_x x + p_y y \le M ,
        \]
        where $p_x,p_y$ are prices and $M$ is income. If $U$ is strictly
        increasing, the budget constraint binds at the optimum.
  \item \textbf{Cost minimisation}: choose
        inputs $(K,L)\ge 0$ to
        \[
          \min\; wL + rK \quad\text{subject to}\quad F(K,L) \ge \bar{Q},
        \]
        producing at least the target output $\bar{Q}$ at input prices
        $w$ (wage) and $r$ (rental rate).
  \item \textbf{Profit maximisation}: choose
        output $q \ge 0$ to
        \[
          \max\; \pi(q) \;=\; R(q) - C(q),
        \]
        revenue minus cost. When the solution is interior ($q^{*}>0$),
        this is effectively unconstrained.
  \item \textbf{Expenditure and portfolio problems}: the
        \emph{expenditure minimisation} problem, $\min\, p_x x + p_y y$ subject to
        $U(x,y) \ge \bar{u}$, is the mirror image (``dual'') of utility
        maximisation; the \emph{portfolio problem} minimises the risk
        (variance) of a wealth allocation subject to the shares summing
        to one and, in richer versions, a target expected return.
\end{itemize}
\end{rulebox}

\begin{rulebox}{Formulating a problem from an economic story}
Work through five questions, in order:
\begin{enumerate}[leftmargin=1.9em,itemsep=0.3ex]
  \item \textbf{Who chooses what?} Name the decision variables and their
        units.
  \item \textbf{What is the goal?} Write the objective function
        explicitly in the decision variables; state max or min.
  \item \textbf{What limits the choice?} List every constraint,
        budgets, technology, time, definitions, nonnegativity.
  \item \textbf{Feasible set.} Collect the constraints into $S$; note
        whether $S$ is compact and which constraints will bind.
  \item \textbf{Classify.} Constrained or (effectively) unconstrained?
        If one equality constraint appears, plan the substitution.
\end{enumerate}
Parameters (prices, income, technology coefficients) are \emph{not}
decision variables: they are held fixed.
\end{rulebox}

\begin{rulebox}{First-order thinking in one variable}
Let $f$ be differentiable. If $x^{*}$ is a local maximum or minimum in
the \emph{interior} of the domain, then
\[
  f'(x^{*}) \;=\; 0
\],
the \textbf{first-order condition} (FOC). Points with $f'(x)=0$ are \emph{critical points}: candidates, necessary but not sufficient
($f(x)=x^{3}$ has $f'(0)=0$ with no extremum). The
\textbf{second-order condition} (SOC) classifies:
$f''(x^{*})<0$ gives a strict local maximum, $f''(x^{*})>0$ a strict
local minimum, $f''(x^{*})=0$ is inconclusive. On a closed interval
$[a,b]$ the global optima exist by Weierstrass and lie among the
\emph{candidates}: interior critical points and the endpoints $a,b$,
evaluate $f$ at each and compare.
\end{rulebox}

\section{Worked Example}

\begin{worked}{Turn a story into a problem and solve it: a firm faces inverse demand $p=120-2q$ and total cost $C(q)=200+20q+\tfrac12q^{2}$}
\wstep{Step 1. Identify the four components before writing any calculus.}
\begin{center}\footnotesize
\begin{tabular}{ll}
  \hline
  component & here\\
  \hline
  choice variable & $q$, the quantity produced\\
  objective & profit $\pi(q)=\text{revenue}-\text{cost}$, to be maximised\\
  constraints & $q\ge0$ (and implicitly $q$ finite)\\
  parameters & the numbers $120$, $2$, $200$, $20$, $\tfrac12$\\
  \hline
\end{tabular}
\end{center}
\wtrap{The price $p$ is \emph{not} a choice variable here --- it is
determined by $q$ through the demand curve. Treating both as free gives two
first-order conditions and an unsolvable or wrong system. Decide what the
firm actually chooses, and substitute everything else out.}

\wstep{Step 2. Build the objective explicitly.}
\[
  R(q)=p\cdot q=(120-2q)q=120q-2q^{2} ,
\]
\[
  \pi(q)=R(q)-C(q)=120q-2q^{2}-\Bigl(200+20q+\tfrac12q^{2}\Bigr)
        =100q-\tfrac52q^{2}-200 .
\]

\wstep{Step 3. First-order condition.}
\[
  \pi'(q)=100-5q=0
  \quad\Longrightarrow\quad
  q^{*}=20 .
\]
Equivalently $MR=MC$: $MR=120-4q$ and $MC=20+q$, equal when
$120-4q=20+q$, i.e.\ $q=20$. \checkmark\ The two routes are the same
equation rearranged, and having both is a free check.

\wstep{Step 4. Second-order condition --- do not skip it.}
\[
  \pi''(q)=-5<0 ,
\]
so the critical point is a maximum, and since $\pi''<0$ \emph{everywhere},
$\pi$ is strictly concave and the maximum is \textbf{global}. That last
sentence is what upgrades a candidate into an answer.
\wtrap{A first-order condition finds points where the function is flat ---
maxima, minima and inflections alike. Reporting $q^{*}=20$ without checking
$\pi''$ is an incomplete solution, and on a cubic cost function it can be
the wrong one.}

\wstep{Step 5. Check the sign constraint.}
$q^{*}=20>0$, so $q\ge0$ is slack and the interior solution stands. Then
\[
  p^{*}=120-40=80,
  \qquad
  \pi^{*}=100(20)-\tfrac52(400)-200=2000-1000-200=800 .
\]

\wstep{Step 6. Should the firm produce at all?}
Profit is $800>0$, so yes. But note what happens if the fixed cost were
$1200$ instead of $200$: $q^{*}$ would be \textbf{unchanged} at $20$, since
a constant does not affect $\pi'$, while profit would become $-200$ and the
firm should shut down.
\wtrap{Fixed costs never move the optimal quantity; they only decide whether
to operate. The first-order condition cannot see them, so the shutdown
question must be asked separately --- comparing the optimised profit against
the alternative of producing nothing.}

\wstep{Step 7. Local, global, existence --- the three questions behind every problem in this volume.}
\begin{itemize}[leftmargin=1.6em,itemsep=0.25ex]
  \item \textbf{Does an optimum exist?} Weierstrass guarantees it on a
        closed bounded feasible set with a continuous objective. Here the
        set $q\ge0$ is unbounded, so existence came instead from
        $\pi\to-\infty$ as $q\to\infty$.
  \item \textbf{Is the candidate a maximum?} The second-order condition.
  \item \textbf{Is it global?} Concavity of the objective on a convex
        feasible set: exactly the content of Chapters~\chnum{optimization}{2} and~\chnum{optimization}{3}.
\end{itemize}

\wstep{Step 8. The same skeleton, constrained.}
A consumer maximising $U(x,y)$ subject to $p_{x}x+p_{y}y\le M$ has the
identical structure: two choice variables instead of one, and an inequality
constraint that will bind (more is better, so the budget is spent). The only
new machinery is the multiplier of \chref{optimization}{7}. Everything else (identify
the choice variables, build the objective, first order, second order, check
the constraints) is what you have just done.
\end{worked}

\section{Practice Problems}

Problems 1--9 drill formulation, classification, and existence; problems
10--20 add the one-variable first-order condition. In every formulation
problem, identify the decision variables, the objective function, the
constraints, and the feasible set, and classify the problem as a
constrained or unconstrained maximisation or minimisation.

\begin{problems}
  \item A consumer has income $M=100$ and buys goods $x$ and $y$ at
        prices $p_x=2$ and $p_y=5$; her utility is $U(x,y)=xy$.
        Formulate her problem. Will the budget constraint bind at the
        optimum?
  \item A price-taking firm sells its output $q$ at the market price
        $p=20$ and produces at cost $C(q)=10+2q^{2}$. Formulate the
        profit-maximisation problem. Is it effectively unconstrained?
  \item A firm must deliver $\bar{Q}=100$ units of output. Its
        technology is $Q=K^{1/2}L^{1/2}$, the wage is $w=4$ and the
        rental rate of capital is $r=1$. Formulate the cost-minimisation
        problem.
  \item Formulate the expenditure-minimisation problem that is dual to
        Problem~1: reach the utility level $\bar{u}=250$ at prices
        $p_x=2$, $p_y=5$ with minimal spending. How do you expect its
        solution to relate to Problem~1?
  \item An investor puts the share $x$ of her wealth into asset $A$
        (return variance $4$) and the remaining share $1-x$ into asset
        $B$ (return variance $1$); the returns are uncorrelated, so the
        portfolio variance is $\sigma^{2}(x)=4x^{2}+(1-x)^{2}$.
        Formulate the minimum-variance problem.
  \item Classify each problem as a constrained or (effectively)
        unconstrained maximisation or minimisation:
        (a) a price taker chooses $q\ge 0$ to maximise $pq-C(q)$,
        expecting $q^{*}>0$;
        (b) $\max\, U(x,y)$ subject to $p_x x+p_y y=M$, $x,y\ge 0$;
        (c) $\min_{x\in\R}\; x^{2}+2x$;
        (d) $\min\, wL+rK$ subject to $F(K,L)=\bar{Q}$, $K,L\ge 0$.
  \item Describe the feasible set
        $B=\{(x,y)\in\R^{2} : 2x+3y\le 12,\; x\ge 0,\; y\ge 0\}$
        geometrically. Is it compact? What does the Weierstrass theorem
        then guarantee for a continuous utility function on $B$?
  \item Which of the following problems possess a solution, and why (or
        why not)?
        (a) $\max\, x$ subject to $0<x<1$;
        (b) $\max\, x$ subject to $0\le x\le 1$;
        (c) $\min\, 1/x$ subject to $0<x\le 1$;
        (d) $\max\, 1/x$ subject to $0<x\le 1$.
  \item A firm's cost of running $q$ machines is $C(q)=q^{2}-6q+10$.
        Rewrite the minimisation of $C$ as an equivalent maximisation
        problem, solve both, and confirm that they have the same
        solution $q^{*}$ and that the optimal values differ only in
        sign.
  \item Find and classify all critical points of $f(x)=x^{3}-3x$ on
        $\R$. Are the local optima global? Why or why not?
  \item A monopolist faces inverse demand $p(q)=20-q$ and cost
        $c(q)=q^{2}+4q$. Formulate the profit-maximisation problem,
        solve the FOC, verify the SOC, and compute the maximal profit.
  \item A firm faces inverse demand $p(q)=100-q$ and a constant unit
        cost of $5$. Find the profit-maximising output and the maximal
        profit.
  \item A concert promoter faces inverse demand $p(q)=30-0.5q$ for
        tickets and, the hall being already paid for, wants to maximise
        \emph{revenue}. Find the revenue-maximising quantity, the
        ticket price, and the maximal revenue.
  \item A plant has cost function $C(q)=100+q^{2}$ for $q>0$. Find the
        output that minimises \emph{average} cost $AC(q)=C(q)/q$,
        verify the SOC, and confirm that marginal cost equals average
        cost at the minimum.
  \item Solve the cost-minimisation problem of Problem~3 by
        substitution: minimise $4L+K$ subject to $K^{1/2}L^{1/2}=100$.
  \item A consumer with utility $U(x,y)=xy$ faces the budget
        $2x+y=8$. Solve her problem by substituting the budget into the
        objective, and verify the SOC.
  \item Find the global maximum and global minimum of
        $f(x)=x^{3}-12x$ on the interval $[-3,5]$. Which optimum lies
        at a boundary point?
  \item A retailer sells $D=200$ units per year. Each order placed
        costs a fixed fee of $25$; holding inventory costs $4$ per unit
        per year, and average inventory is half the order size $q$.
        Total annual cost is $C(q)=\dfrac{200}{q}\cdot 25+\dfrac{4q}{2}$
        for $q>0$. Formulate and solve the cost-minimisation problem.
  \item A gardener has $40$ metres of fence and builds a rectangular
        enclosure against an existing wall, so only three sides need
        fencing. Formulate the area-maximisation problem in one
        variable and solve it.
  \item A monopolist faces inverse demand $p(q)=12-q$ and produces at a
        unit cost of $2$; the government levies a per-unit tax $t=2$.
        Formulate the profit-maximisation problem and solve it. What
        price do consumers pay?
\end{problems}

\section{Solutions}

\begin{solutions}
  \item \textbf{Decision variables:} the quantities $x,y\ge 0$.
        \textbf{Objective:} maximise $U(x,y)=xy$.
        \textbf{Constraints:} the budget $2x+5y\le 100$ and $x,y\ge 0$.
        \textbf{Feasible set:} the budget set
        $S=\{(x,y): 2x+5y\le 100,\ x\ge 0,\ y\ge 0\}$, a closed and
        bounded triangle. A \emph{constrained maximisation}. Since $U$
        is strictly increasing in both goods, no bundle with unspent
        income can be optimal: the budget binds, $2x+5y=100$.
        (Preview of the solution by substitution: $y=(100-2x)/5$ gives
        $u(x)=x(100-2x)/5$, $u'(x)=(100-4x)/5=0$, so $x^{*}=25$,
        $y^{*}=10$, $U^{*}=250$.)

  \item \textbf{Decision variable:} output $q\ge 0$.
        \textbf{Objective:} maximise profit
        $\pi(q)=20q-(10+2q^{2})$.
        \textbf{Constraint:} $q\ge 0$ only, so the feasible set is
        $[0,\infty)$. A \emph{maximisation}, and effectively
        \emph{unconstrained}: the FOC $\pi'(q)=20-4q=0$ gives
        $q^{*}=5>0$, so the nonnegativity constraint is slack.
        ($\pi''=-4<0$; $\pi(5)=100-10-50=40$.)

  \item \textbf{Decision variables:} inputs $K,L\ge 0$.
        \textbf{Objective:} minimise cost $C=wL+rK=4L+K$.
        \textbf{Constraint:} the output requirement
        $K^{1/2}L^{1/2}\ge 100$ (binding at the optimum, since cost is
        increasing in both inputs: $K^{1/2}L^{1/2}=100$), plus
        $K,L\ge 0$.
        \textbf{Feasible set:}
        $S=\{(K,L): K^{1/2}L^{1/2}\ge 100,\ K,L\ge 0\}$, the region on
        or above the isoquant. A \emph{constrained minimisation}.

  \item \textbf{Decision variables:} $x,y\ge 0$.
        \textbf{Objective:} minimise expenditure $E=2x+5y$.
        \textbf{Constraint:} the utility requirement $xy\ge 250$ (and
        $x,y\ge 0$).
        A \emph{constrained minimisation}: objective and constraint
        have swapped roles relative to Problem~1. By duality one
        expects the same bundle $(25,10)$ with minimal expenditure
        $100$, the income of Problem~1. Check by substitution:
        $y=250/x$ gives $E(x)=2x+1250/x$,
        $E'(x)=2-1250/x^{2}=0 \Rightarrow x^{2}=625 \Rightarrow
        x^{*}=25$, $y^{*}=10$, $E^{*}=50+50=100$. Confirmed.

  \item \textbf{Decision variable:} the portfolio share $x\in[0,1]$
        (the share of $B$ is then $1-x$: the adding-up constraint is
        already substituted in).
        \textbf{Objective:} minimise
        $\sigma^{2}(x)=4x^{2}+(1-x)^{2}$.
        \textbf{Feasible set:} the compact interval $[0,1]$.
        A \emph{constrained minimisation}; by Weierstrass a solution
        exists, and if it is interior the FOC applies. (It is solved in
        Additional Problem~3.)

  \item (a) \emph{Maximisation}; the constraint $q\ge 0$ is slack at
        $q^{*}>0$, so effectively \emph{unconstrained}.
        (b) \emph{Constrained maximisation}: an equality (budget)
        constraint plus sign constraints.
        (c) \emph{Unconstrained minimisation} over all of $\R$ (FOC:
        $2x+2=0$, $x^{*}=-1$, value $-1$).
        (d) \emph{Constrained minimisation}: an equality (technology)
        constraint plus sign constraints.

  \item $B$ is the triangle with vertices $(0,0)$, $(6,0)$ and
        $(0,4)$: the intersection of the half-plane below the budget
        line $2x+3y=12$ with the first quadrant. It is closed (an
        intersection of closed half-planes, boundary included) and
        bounded (it fits in a ball of radius~$10$ around the origin),
        hence \emph{compact}. By Weierstrass, every continuous utility
        function attains a global maximum (and a global minimum)
        somewhere on $B$: existence is guaranteed before any
        calculus is done.

  \item (a) \textbf{No solution.} The supremum is $1$, but no feasible
        $x$ attains it: the set $(0,1)$ is not closed, so Weierstrass
        does not apply.
        (b) \textbf{Solution exists:} $[0,1]$ is compact and $f(x)=x$
        continuous; the maximum is attained at $x^{*}=1$ with value
        $1$.
        (c) \textbf{Solution exists:} $1/x$ is decreasing on $(0,1]$,
        so its minimum is attained at the included endpoint $x^{*}=1$
        with value $1$. (Weierstrass does not apply (the set is not
        closed at $0$) but existence holds anyway: the theorem is
        sufficient, not necessary.)
        (d) \textbf{No solution.} As $x\to 0^{+}$, $1/x\to\infty$: the
        objective is unbounded above on the feasible set.

  \item By the max--min conversion,
        $\min_{q} C(q) = -\max_{q}\bigl(-C(q)\bigr)$ with the same
        optimiser. \emph{Minimisation:} $C'(q)=2q-6=0$ gives $q^{*}=3$;
        $C''=2>0$ (min); $C(3)=9-18+10=1$.
        \emph{Maximisation:} $-C(q)=-q^{2}+6q-10$ has FOC
        $-2q+6=0$, again $q^{*}=3$; second derivative $-2<0$ (max);
        value $-C(3)=-1$. Same solution $q^{*}=3$, and
        $\min C = 1 = -(-1) = -\max(-C)$, as claimed.

  \item FOC: $f'(x)=3x^{2}-3=0 \Rightarrow x=\pm 1$. SOC:
        $f''(x)=6x$. At $x=1$: $f''(1)=6>0$, a strict local minimum
        with $f(1)=1-3=-2$. At $x=-1$: $f''(-1)=-6<0$, a strict local
        maximum with $f(-1)=-1+3=2$. Neither is global: as
        $x\to\infty$, $f(x)\to\infty$ (beating the local max) and as
        $x\to-\infty$, $f(x)\to-\infty$ (beating the local min). The
        domain $\R$ is not compact, and no existence theorem applies.

  \item \textbf{Formulation:} choose $q\ge 0$ to maximise
        \[
          \pi(q)=p(q)\,q-c(q)=(20-q)q-(q^{2}+4q)=16q-2q^{2}.
        \]
        FOC: $\pi'(q)=16-4q=0 \Rightarrow q^{*}=4$. SOC:
        $\pi''(q)=-4<0$, a strict maximum; since $\pi$ is concave, it
        is global. Profit:
        $\pi(4)=16\cdot 4-2\cdot 16=64-32=32$. (Interior: $q^{*}=4>0$,
        so the sign constraint is slack.)

  \item Profit: $\pi(q)=(100-q)q-5q=95q-q^{2}$. FOC:
        $\pi'(q)=95-2q=0 \Rightarrow q^{*}=47.5$. SOC:
        $\pi''(q)=-2<0$, global maximum by concavity. Profit:
        \[
          \pi(47.5)=95\cdot 47.5-47.5^{2}=4512.5-2256.25=2256.25 .
        \]

  \item Revenue: $R(q)=(30-0.5q)q=30q-0.5q^{2}$. FOC:
        $R'(q)=30-q=0 \Rightarrow q^{*}=30$. SOC: $R''(q)=-1<0$,
        maximum. Price: $p^{*}=30-0.5\cdot 30=15$. Revenue:
        $R(30)=900-450=450$ (check: $p^{*}q^{*}=15\cdot 30=450$).

  \item $AC(q)=\dfrac{100+q^{2}}{q}=\dfrac{100}{q}+q$. FOC:
        \[
          AC'(q)=-\frac{100}{q^{2}}+1=0
          \;\Rightarrow\; q^{2}=100 \;\Rightarrow\; q^{*}=10
        \]
        (taking the positive root). SOC:
        $AC''(q)=200/q^{3}$, so $AC''(10)=0.2>0$: a minimum.
        $AC(10)=10+10=20$. Marginal cost: $MC(q)=C'(q)=2q$, so
        $MC(10)=20=AC(10)$: marginal cost crosses average cost at
        its minimum, as it must.

  \item The constraint binds: $K^{1/2}L^{1/2}=100$, i.e.\
        $KL=10000$, so $K=10000/L$. Substituting,
        \[
          C(L)=4L+\frac{10000}{L}, \qquad
          C'(L)=4-\frac{10000}{L^{2}}=0
          \;\Rightarrow\; L^{2}=2500 \;\Rightarrow\; L^{*}=50 .
        \]
        SOC: $C''(L)=20000/L^{3}>0$ for $L>0$: a minimum (and global,
        since $C$ is convex on $L>0$). Then
        $K^{*}=10000/50=200$ and the minimal cost is
        $C=4\cdot 50+200=200+200=400$. Check the constraint:
        $\sqrt{200\cdot 50}=\sqrt{10000}=100$. Correct.

  \item Budget binds ($U$ increasing): $y=8-2x$. Substituting,
        \[
          u(x)=x(8-2x)=8x-2x^{2}, \qquad
          u'(x)=8-4x=0 \;\Rightarrow\; x^{*}=2 .
        \]
        SOC: $u''(x)=-4<0$, a maximum (global by concavity). Then
        $y^{*}=8-4=4$ and $U^{*}=2\cdot 4=8$. Budget check:
        $2\cdot 2+4=8$. Correct.

  \item By Weierstrass a global max and min exist on the compact
        interval $[-3,5]$. Candidates: interior critical points and the
        endpoints. FOC: $f'(x)=3x^{2}-12=0 \Rightarrow x=\pm 2$, both
        interior. Evaluate all four candidates:
        \[
          f(-3)=-27+36=9, \quad f(-2)=-8+24=16,
        \]
        \[
          f(2)=8-24=-16, \quad f(5)=125-60=65 .
        \]
        Global maximum at the \emph{boundary} point $x^{*}=5$ with
        value $65$; global minimum at the interior point $x^{*}=2$ with
        value $-16$. (The interior critical point $x=-2$ is only a
        local maximum: $16<65$.)

  \item \textbf{Formulation:} choose the order size $q>0$ to minimise
        \[
          C(q)=\frac{5000}{q}+2q
        \]
        (ordering cost $200/q$ orders at $25$ each, plus holding cost
        $4\cdot q/2$). Unconstrained on the open domain $q>0$. FOC:
        \[
          C'(q)=-\frac{5000}{q^{2}}+2=0
          \;\Rightarrow\; q^{2}=2500 \;\Rightarrow\; q^{*}=50 .
        \]
        SOC: $C''(q)=10000/q^{3}>0$: minimum, global by convexity on
        $q>0$. Minimal cost: $C(50)=100+100=200$ per year.

  \item Let $x$ be the length of each of the two sides perpendicular
        to the wall; the side parallel to the wall then has length
        $40-2x$. \textbf{Decision variable:} $x\in(0,20)$.
        \textbf{Objective:} maximise the area
        $A(x)=x(40-2x)=40x-2x^{2}$: the fence constraint has been
        substituted in, so the problem is one-variable and effectively
        unconstrained. FOC: $A'(x)=40-4x=0 \Rightarrow x^{*}=10$. SOC:
        $A''(x)=-4<0$, maximum. Sides $10$, $10$ and
        $40-20=20$ metres; area $A=10\cdot 20=200$ square metres.
        Fence check: $10+10+20=40$. Correct.

  \item \textbf{Decision variable:} $q\ge 0$. \textbf{Objective:}
        with the tax the effective unit cost is $2+t=4$, so
        \[
          \pi(q)=(12-q)q-4q=8q-q^{2} .
        \]
        Effectively unconstrained (interior solution). FOC:
        $\pi'(q)=8-2q=0 \Rightarrow q^{*}=4$. SOC: $\pi''(q)=-2<0$,
        maximum. Profit: $\pi(4)=32-16=16$. Consumers pay
        $p^{*}=12-4=8$. (Check: $\pi=8\cdot 4-4\cdot 4=32-16=16$.)
\end{solutions}

\section{Additional Problems}

\begin{problems}
  \item A worker has $16$ hours per day to split between leisure
        $\ell$ and paid work at wage $w=5$; consumption equals labour
        income, $c=5(16-\ell)$, and utility is $U(c,\ell)=c\,\ell$.
        Formulate the problem, reduce it to one variable by
        substitution, and solve it.
  \item Classify each problem as a constrained or (effectively)
        unconstrained maximisation or minimisation, and solve (b) and
        (c): (a) $\min\, x^{2}+y^{2}$ subject to $x+y=1$;
        (b) $\max_{q\in\R}\, -(q-3)^{2}$;
        (c) a firm chooses capital $K\ge 0$ to maximise
        $f(K)=\sqrt{K}-0.1K$, expecting $K^{*}>0$.
  \item Solve the portfolio problem of Practice Problem~5: minimise
        $\sigma^{2}(x)=4x^{2}+(1-x)^{2}$ over $x\in[0,1]$.
  \item A monopolist faces inverse demand $p(q)=50-2q$ and cost
        $C(q)=2q^{2}+10q$. Find the profit-maximising output and the
        maximal profit.
  \item For the cost function $C(q)=48+3q^{2}$, $q>0$, minimise the
        average cost and confirm that $MC=AC$ at the minimum.
  \item A consumer with utility $U(x,y)=xy^{2}$ faces the budget
        $x+y=9$. Solve by substitution and verify the SOC.
  \item Which of the following problems possess a solution?
        (a) $\min\, 1/x$ on $[1,2]$; (b) $\min\, 1/x$ on $(0,\infty)$;
        (c) $\min\, x+1/x$ on $(0,\infty)$. Justify each answer, and
        solve those that have solutions.
  \item Find the global maximum and minimum of $f(x)=x^{2}-4x+1$ on
        $[0,5]$.
  \item A farmer has $60$ metres of fence for a free-standing
        rectangular pen (all four sides fenced). Formulate the
        area-maximisation problem in one variable and solve it.
  \item A firm faces inverse demand $p(q)=24-2q$ and unit cost $4$.
        Find the optimal output, price and profit without a tax, and
        then with a per-unit tax $t=4$. By how much does the consumer
        price rise?
\end{problems}

\section{Additional Solutions}

\begin{solutions}
  \item \textbf{Decision variables:} leisure $\ell\in[0,16]$ and
        consumption $c\ge 0$; \textbf{objective:} maximise
        $U=c\,\ell$; \textbf{constraint:} $c=5(16-\ell)$ (time and
        budget combined). A constrained maximisation with one equality
        constraint. Substituting,
        \[
          u(\ell)=5(16-\ell)\,\ell=80\ell-5\ell^{2}, \qquad
          u'(\ell)=80-10\ell=0 \;\Rightarrow\; \ell^{*}=8 .
        \]
        SOC: $u''(\ell)=-10<0$, maximum. Hence the worker works
        $16-8=8$ hours, consumes $c^{*}=5\cdot 8=40$, and
        $U^{*}=40\cdot 8=320$.

  \item (a) \emph{Constrained minimisation} (equality constraint).
        (b) \emph{Unconstrained maximisation}: FOC $-2(q-3)=0$ gives
        $q^{*}=3$; second derivative $-2<0$; value $0$.
        (c) \emph{Effectively unconstrained maximisation}: FOC
        \[
          f'(K)=\frac{1}{2\sqrt{K}}-0.1=0
          \;\Rightarrow\; \sqrt{K}=5 \;\Rightarrow\; K^{*}=25>0,
        \]
        so $K\ge 0$ is slack. SOC: $f''(K)=-\tfrac14K^{-3/2}<0$,
        maximum. Value $f(25)=5-2.5=2.5$.

  \item FOC:
        $\dd{\sigma^{2}}{x}=8x-2(1-x)=10x-2=0
        \Rightarrow x^{*}=\tfrac15$, which lies in the interior of
        $[0,1]$. SOC: second derivative $10>0$, minimum (global by
        convexity). Minimal variance:
        \[
          \sigma^{2}\bigl(\tfrac15\bigr)
          =4\cdot\tfrac{1}{25}+\bigl(\tfrac45\bigr)^{2}
          =\tfrac{4}{25}+\tfrac{16}{25}=\tfrac{20}{25}=\tfrac45 .
        \]
        Diversification pays: $\tfrac45$ is below the variance $1$ of
        holding asset $B$ alone.

  \item Profit:
        $\pi(q)=(50-2q)q-2q^{2}-10q=50q-2q^{2}-2q^{2}-10q
        =40q-4q^{2}$. FOC: $\pi'(q)=40-8q=0 \Rightarrow q^{*}=5$.
        SOC: $\pi''(q)=-8<0$, global maximum. Profit:
        $\pi(5)=200-100=100$.

  \item $AC(q)=\dfrac{48}{q}+3q$. FOC:
        $AC'(q)=-\dfrac{48}{q^{2}}+3=0 \Rightarrow q^{2}=16
        \Rightarrow q^{*}=4$. SOC: $AC''(q)=96/q^{3}>0$, minimum.
        $AC(4)=12+12=24$. Marginal cost $MC(q)=6q$ gives
        $MC(4)=24=AC(4)$. Confirmed.

  \item Budget binds: $x=9-y$. Substituting,
        \[
          u(y)=(9-y)y^{2}=9y^{2}-y^{3}, \qquad
          u'(y)=18y-3y^{2}=3y(6-y)=0
          \;\Rightarrow\; y=0 \text{ or } y=6 .
        \]
        $y=0$ gives $U=0$ (a minimum on the budget line). At
        $y^{*}=6$: SOC $u''(y)=18-6y$, $u''(6)=-18<0$, maximum. Then
        $x^{*}=3$ and $U^{*}=3\cdot 36=108$. Budget check: $3+6=9$.
        Correct.

  \item (a) \textbf{Yes:} $[1,2]$ is compact and $1/x$ continuous
        there (Weierstrass); since $1/x$ is decreasing, the minimum is
        at $x^{*}=2$ with value $\tfrac12$.
        (b) \textbf{No:} as $x\to\infty$, $1/x\to 0$, so the infimum
        $0$ is never attained, no minimum exists.
        (c) \textbf{Yes}, despite the non-compact domain: FOC
        $1-1/x^{2}=0$ gives $x^{*}=1$ (taking $x>0$); SOC
        $f''(x)=2/x^{3}$, $f''(1)=2>0$, minimum, with value
        $f(1)=1+1=2$. Since $f(x)\to\infty$ both as $x\to 0^{+}$ and as
        $x\to\infty$, the interior critical point is the global
        minimum. Weierstrass is sufficient for existence, not
        necessary.

  \item Candidates: FOC $f'(x)=2x-4=0$ gives the interior point
        $x=2$, plus the endpoints $0$ and $5$. Values:
        \[
          f(0)=1, \qquad f(2)=4-8+1=-3, \qquad f(5)=25-20+1=6 .
        \]
        Global minimum at $x^{*}=2$ with value $-3$; global maximum at
        the boundary point $x^{*}=5$ with value $6$.

  \item With sides $x$ and $y$ and perimeter $2x+2y=60$, substitute
        $y=30-x$: maximise $A(x)=x(30-x)=30x-x^{2}$ for
        $x\in(0,30)$. FOC: $A'(x)=30-2x=0 \Rightarrow x^{*}=15$. SOC:
        $A''(x)=-2<0$, maximum. Hence $y^{*}=15$ (a square) and
        $A=15\cdot 15=225$ square metres. Fence check:
        $2\cdot 15+2\cdot 15=60$. Correct.

  \item \emph{Without tax:} $\pi(q)=(24-2q)q-4q=20q-2q^{2}$. FOC:
        $20-4q=0 \Rightarrow q^{*}=5$; SOC $-4<0$. Price
        $p^{*}=24-10=14$; profit $\pi=100-50=50$.
        \emph{With the tax:} unit cost $4+4=8$, so
        $\pi(q)=(24-2q)q-8q=16q-2q^{2}$. FOC:
        $16-4q=0 \Rightarrow q^{*}=4$; SOC $-4<0$. Price
        $p^{*}=24-8=16$; profit $\pi=64-32=32$. The consumer price
        rises from $14$ to $16$: by $2$, exactly half of the tax:
        with linear demand and constant unit cost the burden is split
        evenly between consumers and the firm.
\end{solutions}
